In this project, we developed a machine learning model to detect trees in urban areas using free Sentinel-2 satellite images. Knowing exactly where trees are is very important for our partner \textit{infrared.city} to simulate how hot cities get in the summer. We combined public tree datasets from different European cities with OpenStreetMap data to train a special dual-head U-Net model. Because the public city data only lists trees on streets and in parks, but completely misses trees in private gardens, our formal evaluation scores looked rather low at first. However, visual checks showed that the model actually works very well and successfully learned to find these missing private trees. By using images from three different seasons (spring, summer, and autumn) and a temporal attention mechanism, the model also learned to handle different climates and changing leaf colors. In the end, we built a working, scalable prototype that can be used in the future to improve temperature and microclimate simulations for cities.
